% !TEX program = xelatex
% AI-effects 프로젝트 — "AI"라는 용어가 지칭해 온 세 가지 역사적 흐름
% (기계화·자동화 / 4차 산업혁명(IoT+AI) / LLM·생성형 AI)에 따라
% reference/ 소장 문헌을 분류하고, 흐름별로 AI의 정의·연구주제·중심변수·
% 내용이 어떻게 변화했는지 정리하는 작업 문헌(report/).
% 01_AI영향_실증연구_정리.tex(실증연구 개관)·02_AI노출도_지표_비교.tex(노출도
% 지표 심화)에 이어, "AI 개념 자체의 역사적 층위"를 축으로 문헌을 재배열한다.
% 컴파일: xelatex -> biber -> xelatex -> xelatex (kotex, biblatex/biber 필요)
\documentclass[11pt]{article}

\usepackage{fontspec}
\usepackage{kotex}
\usepackage[a4paper,margin=2.2cm]{geometry}
\usepackage{longtable}
\usepackage{booktabs}
\usepackage{array}
\usepackage{amsmath}
\usepackage{xcolor}
\usepackage{hyperref}
\usepackage{enumitem}
\usepackage{titlesec}
\usepackage[backend=biber,style=authoryear,maxcitenames=2,uniquelist=false,sorting=nyt]{biblatex}
\addbibresource{../../reference/references.bib}

\hypersetup{colorlinks=true,linkcolor=blue!50!black,citecolor=blue!50!black,urlcolor=blue!50!black}
\newcolumntype{L}[1]{>{\raggedright\arraybackslash}p{#1}}
\setlength{\LTleft}{0pt}
\setlength{\LTright}{0pt}
\renewcommand{\arraystretch}{1.15}
\setlength{\tabcolsep}{4pt}

\title{``AI''는 언제나 같은 뜻이 아니었다:\\[1mm]
자동화 · 4차 산업혁명 · LLM, 세 흐름에 따른 문헌 분류\\[2mm]
\large --- 정의, 연구주제, 중심변수, 내용의 변화 ---}
\author{AI-effects 프로젝트}
\date{\today}

\begin{document}
\maketitle

\begin{abstract}
\noindent 경제학·노동시장 문헌에서 ``AI(인공지능)''는 시기에 따라 서로 다른 대상을 지칭해 왔다. 본 문헌은 이를 (1) 기계에 의한 자동화(automation) 일반, (2) 사물인터넷(IoT)·빅데이터와 결합한 4차 산업혁명기의 좁은 의미의 인공지능(narrow AI/머신러닝), (3) 거대언어모형(LLM)·생성형 AI의 등장이라는 세 가지 역사적 흐름으로 구분하고, \texttt{reference/notes/}에 수록된 문헌 54편 전체를 이 세 흐름에 따라 분류한다. 각 흐름에서 ``AI''가 실제로 무엇을 의미했는지(정의), 그로부터 어떤 질문이 제기되었는지(연구주제), 그 질문에 답하기 위해 어떤 변수가 사용되었는지(중심변수), 그리고 방법론과 데이터가 어떻게 달라졌는지(내용)를 정리한다. 특히 세 흐름의 경계에 위치한 문헌들 — 자동화 이론을 생성형 AI에 그대로 적용한 연구, 하나의 논문 안에서 노출도(exposure)와 실사용(usage) 지표를 나란히 쓴 연구, 이론적 노출과 실제 사용 데이터를 하나의 지표로 결합한 연구 — 을 별도로 짚어, 본 프로젝트가 지향하는 ``exposure--usage 결합 분석''이 정확히 어느 지점의 방법론적 공백을 메우려는 시도인지를 논의한다.
\end{abstract}

\tableofcontents

% ============================================================
\section{서론: 왜 ``AI''를 시대별로 구분해야 하는가}
% ============================================================

\texttt{reference/notes/}에는 로봇의 노동대체를 다루는 1980~2010년대 자료 분석부터, 2026년 Claude 대화 로그 100만 건을 분석한 최신 보고서까지 발행 시점이 40여 년에 걸쳐 있는 문헌 54편이 모여 있다. 이 문헌들은 모두 ``AI가 노동시장·거시경제에 미치는 영향''을 다룬다고 말할 수 있지만, 정작 ``AI''라는 단어로 지칭하는 대상은 논문마다 크게 다르다. 어떤 논문에서 AI는 산업용 로봇과 사실상 동의어이고, 어떤 논문에서는 특허 텍스트로 식별되는 머신러닝 응용 일반을 뜻하며, 또 다른 논문에서는 ChatGPT·Claude 같은 특정 생성형 AI 제품의 실제 대화 로그를 뜻한다. 이 차이를 구분하지 않고 문헌을 병렬적으로 인용하면, ``AI 노출도가 고용에 미치는 영향''과 같은 명제가 실제로는 서로 다른 세 가지 기술을 가리키는 세 가지 별개의 명제를 뭉뚱그린 것일 수 있다.

이 문헌은 다음과 같은 세 흐름 구분을 축으로 삼는다.

\begin{enumerate}[label=(\arabic*)]
  \item \textbf{기계화·자동화(Automation)}: 산업혁명 이래 자본이 인간 노동을 물리적·기계적으로 대체해 온 가장 넓고 오래된 흐름. 여기서 ``AI''는 흔히 자동화 기술 일반의 최신 사례로 취급되며, 로봇·전용기계·소프트웨어와 개념적으로 구분되지 않는다.
  \item \textbf{4차 산업혁명(IoT+AI, narrow AI)}: 대략 2010년대 중반부터 2022년 이전까지, 사물인터넷·빅데이터·클라우드와 결합한 머신러닝(딥러닝·컴�터비전·자연어처리 등 특정 응용기술)이 ``AI''로 명명되며 별도의 정책·통계 범주로 자리잡은 흐름. 기업의 AI 도입 여부를 묻는 정부 실태조사, 전문가 크라우드소싱·특허텍스트 기반 직업별 노출지수가 이 시기의 전형적 산물이다.
  \item \textbf{LLM·생성형 AI(2022년 이후)}: ChatGPT(2022.11)·Claude 등 거대언어모형 기반 생성형 AI의 대중적 확산 이후, ``AI''는 구체적인 제품과 그 실제 사용 로그로 연구 가능한 대상이 된다. 노출도 측정 자체가 LLM에게 직접 과업 자동화 가능성을 묻는 방식으로 바뀌고, Anthropic·OpenAI 등이 공개하는 실사용(usage) 데이터가 새로운 실증 기반으로 등장한다.
\end{enumerate}

세 흐름은 시기적으로 이어지지만 서로를 대체하지 않는다 — 2020년대에도 로봇 자동화 이론(흐름 1)은 계속 발전하고 있고, 4차 산업혁명형 기업 AI도입 조사(흐름 2)도 매년 갱신된다. 오히려 최근 문헌일수록 세 흐름의 개념·방법론을 의식적으로 섞어 쓰는 경향이 뚜렷하다(\S\ref{sec:boundary} 참조). 표\,\ref{tab:overview}은 세 흐름을 한눈에 비교한 것이며, \S\ref{sec:era1}--\S\ref{sec:era3}에서 각각을 상세히 다룬다. 부록 표\,\ref{tab:appendix}는 \texttt{reference/notes/} 문헌 54편 전체를 이 분류에 따라 정리한 일람표다.

\begin{longtable}{L{1.8cm}L{2.3cm}L{3.6cm}L{3.6cm}L{3.6cm}}
\caption{세 흐름 개관: AI 정의·연구주제·중심변수·방법론의 변화}\label{tab:overview}\\
\toprule
\textbf{흐름} & \textbf{대략적 시기} & \textbf{AI의 정의} & \textbf{연구주제} & \textbf{중심변수·방법론} \\
\midrule
\endfirsthead
\toprule
\textbf{흐름} & \textbf{대략적 시기} & \textbf{AI의 정의} & \textbf{연구주제} & \textbf{중심변수·방법론} \\
\midrule
\endhead
\midrule \multicolumn{5}{r}{\textit{다음 쪽에 계속}} \\
\endfoot
\bottomrule
\endlastfoot
(1) 기계화·자동화 &
1960년대~현재 (이론), 1990~2010년대 (로봇 실증) &
자본이 인간 과업을 수행할 수 있게 되는 것 일반. 로봇·전용기계·소프트웨어·AI가 개념적으로 구분되지 않음 &
과업대체(displacement)와 신규과업 창출(reinstatement)의 균형, 노동소득분배율·임금불평등의 장기 추이 &
과업변위 지수, 산업용 로봇 보급도, 루틴 직무 비중, 노동소득분배율, 균형성장경로(BGP) 이론모형 \\
(2) 4차 산업혁명 (IoT+AI) &
약 2015~2022년 &
머신러닝·딥러닝·컴퓨터비전 등 특정 알고리즘 응용기술("narrow AI"). 정책·통계상 별도 범주로 취급 &
직업·산업별 AI 노출도와 고용·임금의 관계, 기업의 AI 도입 결정요인과 생산성·경영성과 효과 &
전문가 크라우드소싱 노출지수(AIOE), 특허텍스트 노출지수(Webb), 기업 AI도입 여부 더미(정부 실태조사), TFP \\
(3) LLM·생성형 AI &
2022년 11월(ChatGPT)~현재 &
거대언어모형 기반 생성형 AI 제품(ChatGPT, Claude 등) 자체와 그 실제 사용 &
실제 사용 패턴(과업별 대화점유율, 자동화/증강 비중), LLM 자체평가 기반 노출도, 생산성 이득의 거시적 집계, 노출과 실사용의 결합 &
Anthropic/OpenAI 등 실사용 로그, LLM 자체평가·델파이 노출지수, 과업단위 시간절감률, effective AI coverage, observed exposure \\
\end{longtable}

% ============================================================
\section{흐름 (1): 기계화·자동화 (Automation)}
\label{sec:era1}
% ============================================================

\subsection{AI의 정의}

이 흐름에서 ``AI''는 별도의 범주가 아니다. 핵심 이론적 틀인 과업기반(task-based) 모형에서 분석 대상은 ``자본이 노동을 대체해 특정 과업을 수행할 수 있게 되었는가''라는 일반적 명제이며, 그 자본이 산업용 로봇이든, 전용기계든, 소프트웨어든, 알고리즘이든 모형상 동일하게 취급된다. \textcite{acemoglu-2018-artificial-intelligence-automation-and}가 제목에 ``Artificial Intelligence''를 명시하면서도 본문에서는 AI를 별도로 정의하지 않고 곧바로 ``자동화(자본이 과업 [N-1,I]를 노동 대신 수행하게 되는 것)''라는 일반 개념으로 치환하는 것이 전형적이다. \textcite{korinek-2024-scenarios-for-the-transition-to}은 이를 극한까지 밀어붙여, AI를 ``컴퓨트(compute)로 매개되는 일반화된 자동화 기술''로 정식화하고 완전자동화(AGI)로의 이행 경로를 이론모형으로 다룬다.

\subsection{연구주제와 중심변수}

이 흐름의 연구주제는 일관되게 \emph{전치효과(displacement effect)}와 \emph{복원효과(reinstatement effect)}의 균형이다 \parencite{acemoglu-2019-automation-and-new-tasks-how}. 자동화가 기존 과업에서 노동을 밀어내는 힘과, 신규 과업 창출이 노동을 새로운 영역으로 복원시키는 힘이 노동소득분배율·임금불평등·고용을 어떻게 결정하는지가 핵심 질문이다. \textcite{acemoglu-2011-skills-tasks-and-technologies-implications}의 숙련-과업 리카도 모형이 이론적 출발점을 제공했고, \textcite{acemoglu-2018-the-race-between-man-and}이 정태·동태 일반균형모형으로, \textcite{acemoglu-2022-tasks-automation-and-the-rise}가 500개 인구통계 집단 단위 실증분석으로 이를 완성한다. \textcite{aghion-2017-artificial-intelligence-and-economic-growth}는 여기에 보몰의 비용질병(Baumol's cost disease) 통찰을 결합해, 완전자동화가 진행되어도 자본분배율이 특정 수준에서 안정될 수 있음을 보인다.

중심변수는 \textbf{과업변위(task displacement) 지수, 산업용 로봇 보급도, 루틴 직무 비중, 노동소득분배율}이다. 실증연구는 로봇 대수·특화소프트웨어 도입 같은 물리적·관측 가능한 자동화 지표를 도구변수로 활용한다(\cite{humlum-2019-robot-adoption-and-labor-market}의 로봇세 일반균형분석, \cite{acemoglu-2022-tasks-automation-and-the-rise}의 2SLS 분석). 정의상 특정 세대의 AI 기술(로봇 대 딥러닝 대 LLM)을 구분하지 않으므로, 이 흐름의 실증 데이터는 대개 1980~2010년대에 국한된다.

\subsection{내용의 변화: 이론에서 정량화로}

\textcite{acemoglu-2011-skills-tasks-and-technologies-implications}·\textcite{acemoglu-2018-the-race-between-man-and}이 순수 이론모형으로 전치·복원효과의 존재를 증명한 뒤, \textcite{acemoglu-2018-artificial-intelligence-automation-and}이 이를 ``so-so technologies(그저그런 기술)''라는 정책적으로 중요한 개념으로 단순화했고, \textcite{acemoglu-2019-automation-and-new-tasks-how}·\textcite{acemoglu-2022-tasks-automation-and-the-rise}가 미국 산업·인구통계 자료로 그 크기를 정량화했다. \textcite{humlum-2019-robot-adoption-and-labor-market}은 이 틀을 로봇이라는 단일 자본재에 특화해 일반균형 동학과 정책(로봇세)의 함의를 도출한다. 가장 최근인 \textcite{korinek-2024-scenarios-for-the-transition-to}은 이 이론적 틀을 완전자동화·AGI 시나리오로 외삽(extrapolate)해, 흐름 (1)의 개념장치가 흐름 (3)이 다루는 첨단 AI의 장기 시나리오 분석에도 그대로 사용될 수 있음을 보여준다 — 즉 자동화 이론은 특정 시대에 갇힌 것이 아니라, 이후 모든 흐름의 공통 문법(grammar)으로 기능한다.

% ============================================================
\section{흐름 (2): 4차 산업혁명 (IoT + AI, narrow AI)}
\label{sec:era2}
% ============================================================

\subsection{AI의 정의}

이 흐름에서 처음으로 ``AI''가 로봇·자동화 일반과 구분되는 독자적 범주로 정의된다. 정의 방식은 크게 두 갈래다. 첫째, \textbf{응용기술 목록형} 정의로, \textcite{felten-2021-occupational-industry-and-geographic-exposure}의 AIOE(AI Occupational Exposure)가 대표적이다 — Electronic Frontier Foundation이 정리한 이미지인식·언어모델링 등 10개 AI 응용을 크라우드소싱으로 O*NET 직업능력과 연결한다. 둘째, \textbf{특허텍스트형} 정의로, \textcite{webb-2020-impact-artificial-intelligence-labor}가 지도학습·강화학습·신경망·딥러닝 관련 특허문서와 직무기술서의 텍스트 중복도로 AI 노출을 식별한다. 이 두 정의 모두 인간(크라우드워커·특허심사관)의 사전 판단에 의존하며, 특정 제품이 아니라 ``머신러닝 응용기술''이라는 추상적 범주를 가리킨다. 정부·산업 통계에서는 AI가 빅데이터·IoT·클라우드와 묶여 ``4차 산업혁명 기술''로 통칭되기도 한다(\cite{kiet-2021-ai-adoption-productivity}의 BIM 지표).

\subsection{연구주제와 중심변수}

두 갈래의 연구주제가 병행된다. (a) \textbf{직업·산업 단위 노출-고용 관계}: AIOE·Webb 지수가 실제 고용·임금 변화를 예측하는지를 국가별로 검증한다 — 미국 \parencite{acemoglu-2022-artificial-intelligence-and-jobs-evidence}, 한국 \parencite{cheon-2022-ai-employment-wage-effects,han-2023-ai-and-labor-market-change,kiet-2025-ai-occupational-employment-effect,kim-2025-youth-job-choice-ai-exposure}. (b) \textbf{기업 단위 AI 도입 결정요인·성과효과}: 정부 실태조사(기업활동조사·사업체패널조사)로 AI 도입 여부와 생산성·매출·고용의 관계를 추정한다 — 한국 \parencite{kim-2021-ai-adoption-status-major-industries,kiet-2024-ai-adoption-firms-policy,kiet-2024-ai-labor-market-industry,kim-2025-ai-adoption-firm-performance,lee-2025-ai-adoption-determinants-performance,spri-2024-ai-industry-survey}, 미국 \parencite{hampole-2025-artificial-intelligence-and-the-labor}(이력서 텍스트로 기업의 실제 AI응용 도입을 식별).

중심변수는 \textbf{직업/산업 단위 AI 노출지수(AIOE, Webb score), 기업 단위 AI 도입 여부 더미, 총요소생산성(TFP), 매출·부가가치}다. 이 흐름의 실증연구가 공유하는 가장 뚜렷한 특징은, 노출지수가 여전히 \emph{ex-ante(사전적)} 지표라는 점이다 — 실제로 AI가 사용되었는지가 아니라, 기술적으로 사용 가능한지를 전문가·특허·서베이로 판단한다. \textcite{chang-2022-deep-learning-job-classification-system}처럼 이 시기 AI(딥러닝·BERT류 언어모형)가 노출의 \emph{대상}이 아니라 연구자가 구인공고를 분류하는 \emph{도구}로 쓰이는 사례도 있다 — 이는 4차 산업혁명기 머신러닝이 이미 연구 인프라로 일상화되었음을 보여준다.

\subsection{내용의 변화: 미국 방법론의 이식과 기업조사의 반복}

\textcite{felten-2021-occupational-industry-and-geographic-exposure}·\textcite{webb-2020-impact-artificial-intelligence-labor}가 미국에서 AIOE·특허기반 지수를 구축한 직후, 한국 문헌들은 이 방법론을 그대로 이식해 한국 자료에 적용한다 — \textcite{cheon-2022-ai-employment-wage-effects}(AIOE, 고용보험DB), \textcite{han-2023-ai-and-labor-market-change}(Webb 지수, 로봇·소프트웨어와 비교). 동시에 한국 정부기관(KIET·KISDI·KLI·SPRI)은 매년 기업활동조사·사업체패널조사에 AI 도입 문항을 추가해 시계열 실태조사를 축적한다 — \textcite{kiet-2021-ai-adoption-productivity}(2021, IV추정)부터 \textcite{lee-2025-ai-adoption-determinants-performance}(2025, 다시점 DID)까지 방법론이 인과추정 쪽으로 정교화되지만, 도입률 자체는 5\%대에서 좀처럼 오르지 않는다는 관찰이 반복된다(\cite{spri-2024-ai-industry-survey}: 4.5\%$\to$6.4\%). 흐름 (2)의 문헌들은 대체로 2019~2022년 자료(혹은 그 이전)에 기반하며, 여러 논문이 ``분석기간이 생성형 AI 확산 이전이라 결과 해석에 유의해야 한다''고 스스로 명시한다(\cite{kiet-2025-ai-occupational-employment-effect}; \cite{hampole-2025-artificial-intelligence-and-the-labor}) — 이는 흐름 (2)에서 (3)으로의 전환이 데이터 상의 단절로 나타남을 시사한다.

% ============================================================
\section{흐름 (3): LLM · 생성형 AI}
\label{sec:era3}
% ============================================================

\subsection{AI의 정의}

2022년 11월 ChatGPT 공개 이후, ``AI''는 다시 한번 좁아진다 — 이제는 머신러닝 일반이 아니라 \textbf{거대언어모형(LLM) 기반 생성형 AI}라는 구체적 제품군(ChatGPT, Claude, Gemini 등)을 가리킨다. 이 흐름은 노출도 측정 방식 자체를 바꾼다. \textcite{eloundou-2023-gpts-are-gpts-an-early}가 처음으로 ``LLM(GPT-4) 자신에게 과업의 자동화 가능성을 판정하게 하는'' 방법을 제시했고, 한국의 \textcite{son-2025-llm-based-ai-job-exposure-measurement}·\textcite{keis-2025-ai-job-exposure-labor-demand}는 복수의 오픈소스 LLM을 ``전문가 패널''로 삼는 델파이 앙상블로 이를 확장한다. 동시에 이 흐름에서 처음으로 ``AI''가 \emph{측정 대상}이 아니라 \emph{실제 관측되는 사용 로그}로 존재하게 된다 — Anthropic Economic Index(AEI) 시리즈 \parencite{handa-2025-which-economic-tasks-are,appel-2025-uneven-geographic-and-enterprise-ai,appel-2026-economic-primitives,massenkoff-2026-cadences,massenkoff-2026-learning-curves,tamkin-2025-estimating-ai-productivity-gains-from}가 Claude 대화 수백만 건을, Real-Time Population Survey(RPS) 계열 \parencite{bick-2026-mind-the-gap-ai-adoption,bick-2026-what-work-does-generative}과 한국은행 가계조사 \parencite{bok-2025-rapid-ai-diffusion-productivity-effects,bok-2026-ai-adoption-productivity-effects}가 근로자 서베이 응답을 직접 자료로 쓴다.

\subsection{연구주제와 중심변수}

연구주제가 뚜렷하게 세 갈래로 확장된다. (a) \textbf{실사용 패턴 자체의 기술(記述)}: 어떤 과업에 AI가 실제로 쓰이는지, 자동화(위임)인지 증강(협업)인지를 대화 로그로 직접 관측한다 \parencite{handa-2025-which-economic-tasks-are,appel-2025-uneven-geographic-and-enterprise-ai}. (b) \textbf{생산성 이득의 거시적 집계}: 과업 단위 시간절감률을 Hulten 정리로 경제 전체 TFP·GDP 효과로 환산한다 \parencite{acemoglu-2024-the-simple-macroeconomics-of-ai,tamkin-2025-estimating-ai-productivity-gains-from,bok-2025-rapid-ai-diffusion-productivity-effects,nam-2026-ai-macroeconomic-impact-analysis}. (c) \textbf{채택(adoption)의 결정요인과 노동시장 초기효과}: 무엇이 개인·기업의 생성형 AI 채택을 좌우하는지, 청년 고용·연공서열에 어떤 초기 징후가 나타나는지를 다룬다 \parencite{bick-2026-what-work-does-generative,bick-2026-mind-the-gap-ai-adoption,han-2025-ai-diffusion-youth-employment-decline,bok-2026-youth-employment-decline-ai-career-ladder,kim-2025-youth-job-choice-ai-exposure,min-2025-productivity-gains-generative-ai,noh-2025-ai-based-manufacturing-innovation-employment}.

중심변수도 흐름 (2)와 질적으로 달라진다. 단순한 노출 더미 대신 \textbf{과업별 실제 사용/채택률(자기보고 또는 대화 로그 관측), 자동화 대 증강 비중, 과업 성공률(task success rate), 효과적 AI 커버리지(effective AI coverage), AI 자율성 점수}처럼 ``실제로 무슨 일이 일어나고 있는가''를 담은 다차원 지표가 쓰인다(\cite{appel-2026-economic-primitives}의 5개 신규 ``경제 원시지표''가 대표적). 노출지수 자체도 진화한다 — \textcite{felten-2023-how-will-language-modelers-like-chatgpt}는 자신이 2021년에 만든 흐름 (2)형 AIOE를 ``언어모델링'' 응용에만 재가중해 흐름 (3)으로 직접 이행하는 가교 역할을 한다.

\subsection{내용의 변화: 이론적 노출에서 관측된 사용으로}

흐름 (2)의 핵심 긴장 — 노출지수가 예측하는 것과 실제 일어나는 일이 얼마나 다른가 — 이 흐름 (3)에서 정면으로 검증된다. \textcite{bick-2026-what-work-does-generative}는 6개의 노출점수(Eloundou α/β/ζ, Webb, Felten, Brynjolfsson-Mitchell-Rock)가 실제 채택률 변이의 절반 이하만 설명함을 보이고, 챗로그 기반 과업점유율(Anthropic·OpenAI·Microsoft)이 서베이 기반 실제 채택 패턴과 상관 0.4 미만으로 낮다는 것도 함께 보인다. 이는 ``노출 가능성''과 ``실제 사용''이 별개의 정보를 담고 있음을 뜻하며, 본 프로젝트가 이 둘을 결합하려는 이유이기도 하다(\S\ref{sec:position}).

한국 문헌에서는 한국노동연구원(KLI)의 장지연 계열 연구 \parencite{chang-2024-ai-development-employment-effects,chang-2025-ai-era-skills,chang-2026-ai-technology-diffusion-employment,chang-2026-skill-network-ai-era}와 한국은행 조사국의 서동현·오삼일 계열 연구 \parencite{bok-2025-rapid-ai-diffusion-productivity-effects,bok-2026-ai-adoption-productivity-effects,han-2025-ai-diffusion-youth-employment-decline,bok-2026-youth-employment-decline-ai-career-ladder}가 각각 GPT-4 델파이 숙련노출도와 국내 최초 대표표본 가계조사를 축으로 흐름 (2)에서 흐름 (3)으로의 이행을 국내 자료로 반복 검증하고 있다. 특히 한국은행 계열은 \textcite{bok-2025-rapid-ai-diffusion-productivity-effects}(활용률·잠재적 생산성)에서 \textcite{bok-2026-ai-adoption-productivity-effects}(실현된 생산성과의 단절)로 이어지며, ``시간절감이 실제 생산성 증가로 이어지지 않는다''는 흐름 (3) 특유의 발견(생산성 단절, productivity disconnect)을 제시한다 — 이는 흐름 (1)~(2)의 이론이 암묵적으로 가정해 온 ``비용절감=생산성 증가''라는 등식이 생성형 AI 단계에서는 자명하지 않을 수 있음을 시사한다.

% ============================================================
\section{흐름 사이의 경계 문헌: exposure와 usage가 만나는 지점}
\label{sec:boundary}
% ============================================================

일부 문헌은 하나의 흐름에 깔끔히 속하지 않고, 두 흐름의 개념·데이터를 의식적으로 병행한다. 이 경계 문헌들이 본 프로젝트의 문제의식과 가장 직접적으로 맞닿아 있다.

\begin{itemize}[leftmargin=*]
  \item \textbf{흐름 (1)$\to$(3) 직접 도약}: \textcite{acemoglu-2024-the-simple-macroeconomics-of-ai}는 흐름 (1)의 과업기반 이론·Hulten 정리를 그대로 유지한 채, 과업 노출 비중만 \textcite{eloundou-2023-gpts-are-gpts-an-early}의 GPT-4 기반 측정으로 교체한다. 자동화 이론이 특정 기술 세대에 종속되지 않는 범용 문법임을 보여주는 가장 명확한 사례다.
  \item \textbf{흐름 (2)$\to$(3) 재가중}: \textcite{felten-2023-how-will-language-modelers-like-chatgpt}는 자신의 2021년 AIOE(흐름 2) 구성요소 중 ``언어모델링'' 응용만 남기고 재계산해 LLM 특화 지수로 전환한다.
  \item \textbf{한 논문 안에서 exposure와 adoption/usage를 나란히 배치}: \textcite{chang-2024-ai-development-employment-effects}는 AIOE(exposure, 흐름 2)와 GPT Exposure(흐름 3) 및 기업 AI 도입 여부(adoption)를 명시적으로 구분해 병렬 분석한다. \textcite{noh-2025-ai-labor-coexistence}는 AIOE(2019--2024 고용보험DB)와 생성형 AI 확산기(2023--2025)의 산업 사례 인터뷰를 함께 쓴다. \textcite{seo-2025-ai-adoption-characteristics-business-reports}는 2010--2024년 기업 사업보고서 텍스트마이닝으로 AI 키워드 자체가 머신러닝에서 생성형 AI로 무게중심이 이동하는 과정을 시계열로 직접 포착한다. \textcite{nam-2026-ai-macroeconomic-impact-analysis}는 흐름 (1)의 Acemoglu-Hulten 과업모형, 흐름 (3)식으로 확장된 5축 과업지표, 흐름 (2)식 기업 AI/4차산업혁명기술 활용 실증을 한 논문 안에서 모두 결합한다.
  \item \textbf{이론적 노출과 실제 사용의 명시적 결합(observed exposure)}: \textcite{massenkoff-2026-labor-market-impacts-of-ai}는 \textcite{eloundou-2023-gpts-are-gpts-an-early}의 이론적 노출도($\beta$)를 Anthropic의 실사용 데이터로 게이팅(gating)해 ``관측노출(observed exposure)''이라는 단일 지표로 결합한다 — 기술적으로 가능한 것과 실제로 쓰이는 것의 곱집합만을 ``노출''로 인정하는 접근이다. 이는 exposure 지표와 usage 지표를 통계적으로 연동하려는 본 프로젝트의 목표와 방법론적으로 가장 근접한 선례다.
  \item \textbf{방법론 자체의 계보를 정리하는 메타문헌}: \textcite{lee-2024-generative-ai-labor-market-measurement}은 전문가서베이(흐름 2) $\to$ 특허텍스트(흐름 2) $\to$ LLM 자체평가(흐름 3) $\to$ LLM 델파이(흐름 3)로 이어지는 측정방법론 4계열을 비판적으로 비교해, 이 문헌이 채택하는 3흐름 구분과 사실상 동일한 시대구분을 방법론 축으로 제시한다.
\end{itemize}

% ============================================================
\section{본 프로젝트의 위치}
\label{sec:position}
% ============================================================

\texttt{CLAUDE.md}가 명시하는 이 저장소의 목표 — exposure 지표(사전적·이론적 노출)와 usage 데이터(실현된 실제 사용)를 결합해 분석하는 것 — 는 \S\ref{sec:boundary}에서 정리한 경계 문헌들, 그중에서도 \textcite{massenkoff-2026-labor-market-impacts-of-ai}의 관측노출 접근과 정확히 같은 문제의식을 공유한다. 세 흐름 구분에 비추어 보면, 이 목표는 ``흐름 (2)식 사전적 노출지수만으로는 실제 무슨 일이 일어나는지 알 수 없고, 흐름 (3)식 실사용 데이터만으로는 아직 AI가 쓰이지 않는 영역에서 무엇이 일어날지 예측할 수 없다''는 두 가지 한계를 동시에 보완하려는 시도로 재서술할 수 있다. \textcite{bick-2026-what-work-does-generative}가 정량적으로 보였듯 노출지수와 실제 채택률의 상관이 낮다는 사실은, 두 자료를 굳이 결합할 필요가 없다는 뜻이 아니라 — 오히려 둘이 서로 다른 정보를 담고 있어 결합할 \emph{가치}가 있다는 뜻으로 해석하는 것이 본 프로젝트의 입장에 부합한다.

% ============================================================
\section{결론}
% ============================================================

``AI가 노동시장에 미치는 영향''이라는 하나의 질문 아래, 실제로는 로봇 자동화 이론(흐름 1), 머신러닝 응용의 직업별 노출과 기업 도입(흐름 2), 생성형 AI의 실제 사용과 채택(흐름 3)이라는 세 개의 서로 다른 — 그러나 이론적으로 연결된 — 연구 프로그램이 존재한다. 흐름이 바뀔 때마다 AI의 정의(자동화 일반 $\to$ 응용기술 목록 $\to$ 구체적 제품), 중심변수(과업변위·노동소득분배율 $\to$ 노출지수·도입 더미 $\to$ 사용률·자동화-증강 비중), 데이터의 성격(관측된 자본재 $\to$ 서베이·특허 $\to$ 실제 대화 로그)이 체계적으로 달라졌다. 그러나 흐름 (1)의 과업기반 이론이 흐름 (3)의 실증에도 그대로 적용되는 사례(\cite{acemoglu-2024-the-simple-macroeconomics-of-ai})가 보여주듯, 세 흐름은 단절이 아니라 같은 질문에 대한 누적적 정교화로 이해하는 것이 타당하다. 본 프로젝트가 구축할 \texttt{data/proc/merged/}는 이 누적적 정교화의 다음 단계 — 흐름 (2)의 노출도와 흐름 (3)의 실사용을 실제로 연동한 실증 데이터셋 — 를 지향한다.

% ============================================================
\appendix
\section{부록: 문헌 54편 분류 일람}
% ============================================================

아래 표는 \texttt{reference/notes/}에 수록된 문헌 전체(54편)를 흐름별로 정리한 것이다. ``경계''로 표시된 문헌은 \S\ref{sec:boundary}에서 논의한 것처럼 둘 이상의 흐름의 개념·데이터를 명시적으로 병행한다.

\begin{longtable}{L{0.9cm}L{4.6cm}L{0.9cm}L{6.2cm}}
\caption{문헌 54편의 흐름별 분류}\label{tab:appendix}\\
\toprule
\textbf{연도} & \textbf{문헌} & \textbf{흐름} & \textbf{비고} \\
\midrule
\endfirsthead
\toprule
\textbf{연도} & \textbf{문헌} & \textbf{흐름} & \textbf{비고} \\
\midrule
\endhead
\midrule \multicolumn{4}{r}{\textit{다음 쪽에 계속}} \\
\endfoot
\bottomrule
\endlastfoot
2011 & \textcite{acemoglu-2011-skills-tasks-and-technologies-implications} & 1 & 숙련-과업 리카도 모형, 후속 자동화 이론의 출발점 \\
2017 & \textcite{aghion-2017-artificial-intelligence-and-economic-growth} & 1 & 보몰 비용질병+자동화모형, 특이점 시나리오 \\
2018 & \textcite{acemoglu-2018-artificial-intelligence-automation-and} & 1 & 제목은 ``AI''이나 내용은 일반 자동화 프레임(so-so technologies) \\
2018 & \textcite{acemoglu-2018-the-race-between-man-and} & 1 & 전치·복원효과 정태·동태 일반균형 이론 \\
2019 & \textcite{acemoglu-2019-automation-and-new-tasks-how} & 1 & 1947--2017 미국 산업자료로 전치·복원효과 정량화 \\
2019 & \textcite{humlum-2019-robot-adoption-and-labor-market} & 1 & 산업로봇 도입의 일반균형 동학, 로봇세 정책분석 \\
2022 & \textcite{acemoglu-2022-tasks-automation-and-the-rise} & 1 & 500개 인구집단 과업변위와 1980--2016 임금구조 \\
2024 & \textcite{korinek-2024-scenarios-for-the-transition-to} & 1 & 완전자동화/AGI 이행 경로 이론모형 \\
2020 & \textcite{webb-2020-impact-artificial-intelligence-labor} & 2 & 경계(1$\leftrightarrow$2): 특허기반 AI노출, 로봇·SW와 병행비교 \\
2021 & \textcite{felten-2021-occupational-industry-and-geographic-exposure} & 2 & AIOE 원 논문, 크라우드소싱 기반 노출지수 \\
2021 & \textcite{kiet-2021-ai-adoption-productivity} & 2 & 제조업 AI도입-생산성 IV 추정 \\
2021 & \textcite{kim-2021-ai-adoption-status-major-industries} & 2 & 5대 산업 AI도입 실태조사(KISDI) \\
2022 & \textcite{acemoglu-2022-artificial-intelligence-and-jobs-evidence} & 2 & Burning Glass 채용공고+Felten/Webb/SML 노출지수(2010--2018) \\
2022 & \textcite{chang-2022-deep-learning-job-classification-system} & 2 & 딥러닝(BERT류)을 분석도구로 사용한 구인공고 분류 \\
2022 & \textcite{cheon-2022-ai-employment-wage-effects} & 2 & AIOE 한국 최초 적용, 고용보험DB \\
2023 & \textcite{han-2023-ai-and-labor-market-change} & 2 & Webb 지수 한국 적용, 로봇·SW와 비교 \\
2024 & \textcite{kiet-2024-ai-adoption-firms-policy} & 2 & 기업활동조사 기반 AI도입 결정요인·생산성(TWFE) \\
2024 & \textcite{kiet-2024-ai-labor-market-industry} & 2 & 기업AI도입+Webb지수, 기업·직종 이중 실증 \\
2025 & \textcite{hampole-2025-artificial-intelligence-and-the-labor} & 2 & 이력서 텍스트 기반 기업 AI응용 식별(2014--2023) \\
2025 & \textcite{kiet-2025-ai-occupational-employment-effect} & 2 & AIOE+고용보험DB 패널(2010--2019), ChatGPT 이전 한계 자인 \\
2025 & \textcite{kim-2025-ai-adoption-firm-performance} & 2 & 기업활동조사(2017--2023), AI도입-TFP 생산성역설 \\
2025 & \textcite{kim-2025-youth-job-choice-ai-exposure} & 2 & Webb기반 AIOE와 청년구직 선택(혼합로짓) \\
2025 & \textcite{lee-2025-ai-adoption-determinants-performance} & 2 & 사업체패널조사, 다시점 DID \\
2025 & \textcite{spri-2024-ai-industry-survey} & 2 & 국가승인통계, AI산업·기업도입률 시계열 \\
2023 & \textcite{eloundou-2023-gpts-are-gpts-an-early} & 3 & LLM 자체평가 노출등급(E0--E3) 방법론 원 논문 \\
2023 & \textcite{felten-2023-how-will-language-modelers-like-chatgpt} & 3 & 경계(2$\to$3): AIOE를 언어모델링 응용으로 재가중 \\
2024 & \textcite{acemoglu-2024-the-simple-macroeconomics-of-ai} & 3 & 경계(1$\to$3): 자동화이론+GPT노출로 거시생산성 추정 \\
2024 & \textcite{chang-2024-ai-development-employment-effects} & 3 & 경계: AIOE·GPT Exposure·기업도입 병행 종합분석 \\
2024 & \textcite{lee-2024-generative-ai-labor-market-measurement} & 3 & 경계: 측정방법론 4계열(2$\to$3) 비판적 메타비교 \\
2025 & \textcite{appel-2025-uneven-geographic-and-enterprise-ai} & 3 & AEI, 지리적 분해+API 기업트래픽 \\
2025 & \textcite{bick-2026-mind-the-gap-ai-adoption} & 3 & (2026) 미국·유럽 7개국 genAI 채택률 서베이 \\
2025 & \textcite{bok-2025-rapid-ai-diffusion-productivity-effects} & 3 & 한국 최초 대표표본 가계조사, 활용률·잠재적 생산성 \\
2025 & \textcite{chang-2025-ai-era-skills} & 3 & GPT-4 델파이 숙련노출도, 숙련 네트워크 구조 \\
2025 & \textcite{handa-2025-which-economic-tasks-are} & 3 & AEI 최초보고서, Claude 실사용 과업매핑 \\
2025 & \textcite{keis-2025-ai-job-exposure-labor-demand} & 3 & LLM 델파이 앙상블 노출도+구인공고 수요변화 \\
2025 & \textcite{min-2025-productivity-gains-generative-ai} & 3 & 생성형AI$\times$분업 상호작용 실험연구 \\
2025 & \textcite{noh-2025-ai-based-manufacturing-innovation-employment} & 3 & GPT-4 델파이+R\&D엔지니어 실제 LLM활용(92.2\%) \\
2025 & \textcite{noh-2025-ai-labor-coexistence} & 3 & 경계(2+3): AIOE(2019--2024)+생성형AI 산업사례 \\
2025 & \textcite{seo-2025-ai-adoption-characteristics-business-reports} & 3 & 경계(2$\to$3): 사업보고서 텍스트마이닝 시계열(2010--2024) \\
2025 & \textcite{son-2025-llm-based-ai-job-exposure-measurement} & 3 & LLM 델파이 노출도 방법론(923개 O*NET 직업) \\
2025 & \textcite{tamkin-2025-estimating-ai-productivity-gains-from} & 3 & Claude 실사용+Hulten정리로 생산성 거시집계 \\
2026 & \textcite{appel-2026-economic-primitives} & 3 & AEI, 5개 신규 경제원시지표(성공률·자율성 등) \\
2026 & \textcite{bick-2026-what-work-does-generative} & 3 & RPS 과업단위 genAI채택지수, 노출점수 설명력 검증 \\
2026 & \textcite{bok-2026-ai-adoption-productivity-effects} & 3 & 시간절감-생산성 단절(disconnect) 실증 \\
2026 & \textcite{bok-2026-youth-employment-decline-ai-career-ladder} & 3 & AIOE+ChatGPT 확산 DDD, 청년고용 위축 \\
2026 & \textcite{chang-2026-ai-technology-diffusion-employment} & 3 & 생성형AI 확산의 직종·기업별 고용효과 \\
2026 & \textcite{chang-2026-skill-network-ai-era} & 3 & 숙련 네트워크 중심성·전이가능성(요약본) \\
2026 & \textcite{fan-2026-aggregate-gains-from-ai} & 3 & IMF, Claude 사용 기반 국가별 AI이득 분배 \\
2026 & \textcite{fan-2026-what-countries-use-ai-and} & 3 & Claude 사용 강도·확산폭의 국가간 결정요인 \\
2026 & \textcite{han-2025-ai-diffusion-youth-employment-decline} & 3 & AIOE+ChatGPT확산 DDD, 청년고용위축(원 선행연구) \\
2026 & \textcite{massenkoff-2026-cadences} & 3 & AEI, 사용 시간대패턴·AI자율성·주관적 인식 \\
2026 & \textcite{massenkoff-2026-labor-market-impacts-of-ai} & 3 & 경계: 이론적노출+실사용 결합 ``관측노출'' 지표 \\
2026 & \textcite{massenkoff-2026-learning-curves} & 3 & AEI, 사용경력(tenure)에 따른 학습곡선 \\
2026 & \textcite{nam-2026-ai-macroeconomic-impact-analysis} & 3 & 경계(1+2+3): 과업모형+5축 지표+기업AI/IR4 실증 \\
\end{longtable}

\clearpage
\printbibliography[heading=bibintoc,title={참고문헌}]

\end{document}
